\subsection{PURE p-VECTORS. GRASSMANNIANS}

Let K be a field and E a vector space over K. A p-vector \(z \in \bigwedge^{p}(\mathbf{E})\) is called pure (or sometimes decomposable) if it is non-zero and there exist vectors \(x_{1}, \ldots, x_{p}\) in E such that \(z = x_{1} \wedge x_{2} \wedge \cdots \wedge x_{p}\) . For this, it is necessary and sufficient that the subspace \(M_{z}\) associated with z (which is always of dimension \(\geqslant p\) for \(z \neq 0\) ) be exactly of dimension p (since \(\bigwedge^{p}(\mathbf{M}_{z})\) is then of dimension 1). In particular, every non-zero scalar, every non-zero element of \(\mathbf{E} = \bigwedge^{1}(\mathbf{E})\) , every non-zero element of \(\bigwedge^{n}(\mathbf{E})\) , when E is of dimension n, is pure.

PROPOSITION 15. Let E be a vector space of dimension n and let e be an element \(\neq0\) of \(\bigwedge^{n}(\mathbf{E})\) (hence forming a basis of this vector space). Let \(\phi:\bigwedge(\mathbf{E})\to\bigwedge(\mathbf{E}^{*})\) be the vector space isomorphism associated with e (no. 11, Proposition 12). If z is a pure element of \(\bigwedge^{p}(\mathbf{E})\) , then \(\phi(z)\) is a pure element of \(\bigwedge^{n-p}(\mathbf{E}^{*})\) and the subspaces associated with z and \(\phi(z)\) are orthogonal.

The cases \(p = 0\) and \(p = n\) are trivial. Suppose therefore that \(1 \leqslant p \leqslant n - 1\) and let \(z = x_1 \wedge \cdots \wedge x_p \neq 0\) . Then there exists a basis \((e_i)_{1 \leqslant i \leqslant n}\) of \(E\) such that \(e_i = x_i\) for \(1 \leqslant i \leqslant p\) and \(e = e_1 \wedge e_2 \wedge \cdots \wedge e_n\) . It then follows from formula (78) of no. 11 that \(\phi(z) = \pm e_{p+1}^* \wedge \cdots \wedge e_n^*\) , whence the proposition.

COROLLARY. If E is of dimension n, every non-zero \((n - 1)\) -vector E is pure.

PROPOSITION 16. For an element \(z \neq 0\) of \(\bigwedge^{p}(\mathbf{E})\) to be pure, it is necessary and sufficient that, for all \(u^{*} \in \bigwedge^{p-1}(\mathbf{E}^{*})\) ,

\[
(u ^ {*} \lrcorner z) \wedge z = 0.\tag{81}
\]

The case \(p = 0\) is trivial and we assume \(p \geqslant 1\) . If \(z = x_1 \wedge \cdots \wedge x_p\) , formula (68) (no. 9) with \(n = p - 1\) shows that \(u^* \lrcorner z\) is a linear combination of the \(x_i\) ( \(1 \leqslant i \leqslant p\) ), whence (81). If on the other hand the subspace \(M_z\) associated with \(z\) is of dimension \(>p\) , consider a basis \((e_j)_{1 \leqslant j \leqslant n}\) of this subspace with \(n >p\) . It follows from no. 11, Proposition 13 that each of the \(e_j\) is of the form \(u^* \lrcorner z\) for some \(u^* \in \bigwedge^{p-1}(E^*)\) and relation (81) therefore implies \(e_j \wedge z = 0\) for \(1 \leqslant j \leqslant n\) . It follows that in the expression \(z = \sum_{H} a_H e_H\) (where H runs through the set of subsets of {[}1, n{]} with \(p\) elements) all the coefficients \(a_H\) are zero, whence \(z = 0\) , contrary to the hypothesis.

The criterion of Proposition 16 is equivalent to writing conditions (81) when \(u^*\) runs through a basis of \(\bigwedge^{p-1}(\mathbf{E}^*)\) . In particular, suppose that \(\mathbf{E}\) is of finite dimension \(n\) and let \((e_i)_{1\leq i\leq n}\) be a basis of \(\mathbf{E}\) . Conditions (81) are then equivalent to the conditions

\[
(82-(J, H)) \quad \langle e _ {\mathrm{J}} ^ {*}, (e _ {\mathrm{H}} ^ {*} \lrcorner z) \wedge z \rangle = 0
\]

for all subsets J, H of \([1, n]\) such that \(\text{Card}(J) = p + 1\) and \(\text{Card}(H) = p - 1\) . Now, if I and \(I'\) are two subsets of \([1, n]\) with p elements, formulae (76) of no. 10 and multiplication table (20) of § 7, no. 8 show that

\[
\langle e _ {\mathrm{J}} ^ {*}, \langle e _ {\mathrm{H}} ^ {*} \lrcorner e _ {\mathrm{I}}) \wedge e _ {\mathrm{I} ^ {\prime}} \rangle = 0
\]

unless there exists an \(i \in [1, n]\) such that \(\mathbf{I} - \mathbf{H} = \{i\}\) and \(\mathbf{J} - \mathbf{I}' = \{i\}\) , in which case

\[
\langle e _ {\mathrm{J}} ^ {*}, (e _ {\mathrm{H}} ^ {*} \lrcorner e _ {\mathrm{I}}) \wedge e _ {\mathrm{I} ^ {\prime}} \rangle = (- 1) ^ {(p - 1) (p - 2) / 2} \varepsilon_ {i, \mathrm{J,H}}\tag{83}
\]

where \(\varepsilon_{i,J,H} = \rho_{(i),H}\rho_{(i),I'}\) ; it can then be said that for \(i\in J\cap \complement H\) , \(\varepsilon_{i,J,H}\) is equal to \(+1\) if the number of element of \(J\) which are \(< i\) and the number of elements of \(H\) which are \(< i\) have the same parity, and \(-1\) otherwise.

It follows immediately that if we write \(z = \sum_{\mathbf{I}} a_{\mathbf{I}} e_{\mathbf{I}}\) , where \(\mathbf{I}\) runs through the set of subsets of \([1, n]\) with p elements, relation \((82-(J, H))\) is equivalent to the relation

\[
(84 - (\mathrm{J}, \mathrm{H})) \quad \sum_ {i \in J \cup \complement_ {\mathrm{H}}} \varepsilon_ {i, J, \mathrm{H}} a _ {J - \{i \}} a _ {\mathrm{H} \cup \{i \}} = 0.
\]

Relations (84) are called Grassman's relations; these are therefore necessary and sufficient conditions (when J describes the set of subsets with \(p + 1\) elements and H the set of subsets with \(p - 1\) elements of \([1, n]\) ) for an element \(z \neq 0\) of \(\bigwedge^{p}(\mathbf{E})\) to be pure.

Note that relations (84) are not independent. For example, for \(n = 4\) and \(p = 2\) , Grassmann's relations reduce to the single relation

\[
a _ {1 2} a _ {3 4} - a _ {1 3} a _ {2 4} + a _ {1 4} a _ {2 3} = 0.\tag{85}
\]

Let \(D_p(E)\) be the subset of \(\bigwedge^p(E)\) consisting of the pure \(p\) -vectors; clearly \(D_p(E)\) is saturated with respect to the equivalence relation between \(u\) and \(v\) : ``there exists \(\lambda \in K^*\) such that \(v = \lambda u\)'' and two elements \(u, v\) of \(D_p(E)\) are equivalent under this relation if and only if the subspaces \(M_u\) and \(M_v\) of \(E\) which are associated with them are the same. Therefore we thus obtain a canonical bijection of the set of \(p\) -dimensional vector subspaces of \(E\) onto the image \(G_p(E)\) of \(D_p(E)\) in the projective space \(P(\bigwedge^p(E))\) associated with \(\bigwedge^p(E)\) . The subset \(G_p(E)\) of \(P(\bigwedge^p(E))\) is called the Grassmannian of index \(p\) of the vector space \(E\) . When \(E\) is finite-dimensional and \((e_i)_{1 \leq i \leq n}\) is a basis of \(E\) , the Grassmannian of index \(p\) is the set of points of \(P(\bigwedge^p(E))\) for which a system of homogeneous coordinates \((a_I)\) (relative to the basis \((e_I)\) of \(\bigwedge^p(E)\) ) satisfies Grassmann's relations (84).

When \(\mathbf{E} = \mathbf{K}^n\) , we sometimes write \(\mathbf{G}_{n,p}(\mathbf{K})\) instead of \(\mathbf{G}_p(\mathbf{K}^n)\) , so that \(\mathbf{G}_{n,1}(\mathbf{K}) = \mathbf{P}_{n-1}(\mathbf{K})\) . The mapping \(\mathbf{M} \mapsto \mathbf{M}^0\) , which associates with every \(p\) -dimensional subspace of \(\mathbf{K}^n\) the orthogonal subspace in \(\mathbf{E}^*\) (identified with \(\mathbf{K}^n\) under the choice of basis dual to the canonical basis of \(\mathbf{K}^n\) ) therefore defines a canonical bijection of \(\mathbf{G}_{n,p}(\mathbf{K})\) onto \(\mathbf{G}_{n,n-p}(\mathbf{K})\) ; Proposition 15 shows that this bijection is the restriction to \(\mathbf{G}_{n,p}(\mathbf{K})\) of a canonical isomorphism of the projective space \(\mathbf{P}(\bigwedge^p (\mathbf{K}^n))\) onto the projective space \(\mathbf{P}(\bigwedge^{n-p}(\mathbf{K}^n))\) .

